\begin{abstract}
정확한 자세 추정은 소형 무인 항공기(UAV)의 안정적인 비행을 위한 핵심 요소이나,
적절한 센서 융합 알고리즘의 선택은 정확도, 연산 비용, 구현 복잡도 간의 비자명한 절충을 수반한다.
본 논문은 소형 멀티로터 시스템의 상태 추정 파이프라인을 특성화하는 두 가지 상호 보완적인 연구를 제시한다.

첫 번째 연구에서는 아두이노 마이크로컨트롤러에 연결된 MPU-9250 관성 측정 장치(IMU) 위에
여섯 가지 자세 추정 필터—자이로스코프 적분, 상보 필터, Madgwick 필터, Mahony 필터,
확장 칼만 필터(EKF), 오차 상태 칼만 필터(ESKF)—를 구현한다.
정적 측정 및 지터 분석을 통해 센서 노이즈를 특성화하고,
MuJoCo 기반 시뮬레이션을 통해 정량적 필터 비교를 위한 실제 자세 기준값을 제공한다.

두 번째 연구에서는 CUAV 7 Nano 비행 컨트롤러를 탑재한 SK450 쿼드로터 플랫폼의
하드웨어 구조와 제어 파이프라인을 분석한다.
PX4 및 ArduPilot 오픈소스 자동 조종 프레임워크에 문서화된
계단식 PID(PPID) 제어기와 온보드 EKF 기반 상태 추정기를
Skydio~2 항공 모델을 사용한 MuJoCo 시뮬레이션으로 검증한다.

결과적으로 ESKF가 시뮬레이션에서 가장 높은 자세 추정 정확도를 달성하였으며,
Mahony 필터는 현저히 낮은 연산 비용으로 경쟁력 있는 성능을 제공하였다.
논의 섹션에서는 IMU 필터 비교 연구에서 도출된 필터 선택 기준이
비행 컨트롤러 상태 추정기 설계에 주는 실용적 시사점을 분석하여 두 연구를 연결한다.
\end{abstract}

\hbox

\begin{IEEEkeywords}
IMU; 자세 추정; 센서 융합; EKF; ESKF; Madgwick; Mahony; 계단식 PID; PX4; ArduPilot; UAV; MuJoCo
\end{IEEEkeywords}
